% ==============================================================================
% PLANTILLA MAESTRA - FILOSOFIA DEL DERECHO / UnADM
% ==============================================================================
% Uso:
% 1. Duplica este archivo y renombra la copia por semana/actividad.
% 2. Actualiza documenttitle, documentsubtitle y documentsubject.
% 3. Lee la planeacion de la semana y NO pegues las instrucciones completas en el
%    producto final. Usa las instrucciones solo como lista de cumplimiento.
% 4. Identifica primero el producto solicitado: mapa, cuadro, glosario, tabla,
%    estudio de caso, reflexion, video, etc. El formato pedido manda sobre el
%    gusto personal por escribir ensayo.
% 5. Revisa bibliografia, toma notas, define categorias y redacta con postura.
% 6. Entrega en PDF con la nomenclatura indicada por la planeacion.
%
% Formato institucional aplicado:
% - Fuente tipo Helvetica.
% - Interlineado 1.5.
% - Citas autor-anio con natbib: usar \citep{clave} o \citet{clave}.
% - Referencias en bibliografia BibTeX.
% - Redaccion academica, clara, con analisis propio.
% - Productos visuales elaborados preferentemente en LaTeX/TikZ: tablas,
%   cuadros comparativos, mapas conceptuales/mentales, lineas de tiempo,
%   diagramas de flujo, esquemas, cronologias, matrices y organizadores.
%
% Criterios derivados de retroalimentacion docente:
% - No limitarse a reproducir fuentes: interpretar, conectar y tomar postura.
% - La pregunta guia antes de entregar debe ser: cual es mi posicion frente a lo
%   que expongo y que razones sostienen esa posicion.
% - Evitar redaccion excesivamente general, repeticion de estructuras y tono de
%   resumen automatico. Cada parrafo debe tener una decision propia.
% - Las fuentes no se mencionan como lista decorativa: cada autor debe sostener
%   una definicion, un criterio, una diferencia, una tension o una aplicacion.
% - Toda referencia final debe estar citada dentro del texto y toda cita dentro
%   del texto debe aparecer en referencias. Revisar coincidencia uno a uno.
% - Incluir conclusion personal cuando la actividad lo solicite.
% - En cuadros comparativos: categorias claras, jerarquias, secuencia y sintesis.
% - Revisar diseno solicitado: producto visual/cuadro, conclusiones y referencias.
% - Si la consigna pide mapa conceptual, cuadro, tabla o glosario, construir ese
%   producto. No convertirlo en texto expositivo salvo que la planeacion lo pida.
% - Usar APA 7 edicion en citas y referencias.
% - No presentar las indicaciones de la actividad dentro del documento final.
% - Si se usa IA como apoyo, el resultado debe estar leido, corregido, apropiado
%   y transformado con criterio personal; nunca entregar texto sin procesamiento.
%
% Protocolo contra observaciones de "IA" o "falta de autenticidad":
% - Abrir cada seccion con una idea propia, no con definiciones genericas.
% - Integrar por lo menos una frase de valoracion: "El punto critico es...",
%   "Desde esta perspectiva...", "La consecuencia juridica es...".
% - Pasar de definicion a funcion: que hace el concepto dentro del sistema.
% - Pasar de funcion a consecuencia: que ocurre si se interpreta mal.
% - En conclusiones, no repetir; responder que se comprendio, que se sostiene y
%   que cambia para la interpretacion o practica juridica.
%
% APA 7 operativo para esta materia:
% - La cita puede ser parafraseada o textual. Si la planeacion dice "explica con
%   tus propias palabras", preferir parafrasis con cita autor-anio.
% - En glosarios, colocar la cita junto a la definicion base y despues desarrollar
%   comentario propio. Formula: definicion base + cita -> funcion -> ejemplo o
%   consecuencia.
% - Usar cita textual solo cuando la frase exacta sea necesaria; incluir pagina:
%   \citep[p.~00]{clave}. No llenar todo el trabajo de comillas.
% - En citas de varias fuentes, no acumular autores sin distinguir aportes. Mejor
%   explicar que aporta cada fuente en una frase breve.
% - No usar URLs sueltas dentro del texto. Registrar la fuente en .bib y citarla.
%
% Criterios de redaccion personal (enfoque_personal.txt + guia estrategica):
% - No escribir para llenar espacio; escribir para posicionar una idea con peso.
% - Cada seccion debe responder: que es, como funciona, por que importa,
%   que implica y cual es la postura.
% - Evitar texto generico. Convertir informacion en criterio.
% - Estructura mental recomendada: problema -> sistema -> consecuencia -> postura.
% - Usar un tono formal, directo, analitico y con intencion.
% - Integrar pensamiento de sistemas: funcion, estructura, control, criterio,
%   costo, decision, consecuencia, limites, aplicacion real.
% - Usar primera persona solo con moderacion: "Desde esta perspectiva",
%   "Considero que", "La posicion que sostengo es". Evitar "yo creo".
%
% Checklist antes de entregar:
% [ ] El producto coincide exactamente con la tecnica solicitada en planeacion.
% [ ] El objetivo esta alineado con la planeacion.
% [ ] Las instrucciones no fueron copiadas como contenido del producto.
% [ ] Hay citas en APA autor-anio dentro del texto, ubicadas donde se usa la idea.
% [ ] La bibliografia final coincide con las citas del cuerpo, sin sobrantes.
% [ ] La introduccion abre con una afirmacion fuerte, no con relleno.
% [ ] El desarrollo explica, interpreta y conecta.
% [ ] Cada fuente cumple una funcion analitica, no solo aparece mencionada.
% [ ] Hay ejemplos juridicos o aplicacion al sistema mexicano cuando proceda.
% [ ] La conclusion incluye postura personal argumentada: posicion, razon y efecto.
% [ ] No aparecen instrucciones de la planeacion copiadas en el cuerpo final.
% [ ] Si la actividad pide producto visual, se construyo con TikZ/LaTeX o se
%     justifico por que se uso una captura externa.
% [ ] Si es mapa conceptual, hay jerarquia, palabras de enlace y relaciones.
% [ ] Si es cuadro comparativo, hay categorias comparables y sintesis critica.
% [ ] Si es glosario, cada concepto tiene definicion propia, cita y aplicacion.
% [ ] Si es caso, se identifican hechos, problema, conceptos, aplicacion y postura.
% [ ] Ortografia, acentos, sintaxis y nombres propios revisados.
% [ ] Si se uso IA como apoyo, el texto fue comprendido, revisado y apropiado.
% ==============================================================================

% CREACION DEL DOCUMENTO
\documentclass[
	spanish,
	letterpaper, oneside
]{article}

% ==============================================================================
% INFORMACION DEL DOCUMENTO
% Actualizar en cada actividad.
% Ejemplos:
% \def\documenttitle {Mapa conceptual de la Filosofia del Derecho}
% \def\documentsubtitle {Actividad 1 - Filosofia del Derecho}
% \def\documentsubject {Actividad Semana 2}
% ==============================================================================
\def\documenttitle {Reflexión y estudio de caso: interpretación jurídica}
\def\documentsubtitle {Actividad 5 - Filosofía del Derecho}
\def\documentsubject {Semana 7 - Interpretación jurídica}

\def\documentauthor {Martín Jonathan de la Cruz}
\def\coursename {Filosofía del Derecho}
\def\coursecode {FDE30}

\def\universityname {Universidad Abierta y a Distancia de México}
\def\universityfaculty {Licenciatura en Derecho}
\def\universitydepartment {Filosofía del Derecho}
\def\universitydepartmentimage {departamentos/UnADM}
\def\universitydepartmentimagecfg {height=1.57cm}
\def\universitylocation {Roma Norte, Ciudad de Mexico}

% INTEGRANTES, PROFESORES Y FECHAS
\def\authortable {
	\begin{tabular}{ll}
		\textbf{Alumno:}
		& \begin{tabular}[t]{l}
			 Martín Jonathan de la Cruz \\
		\end{tabular} \\
		Matrícula: & ES2611202040 \\

		\textit{Profesor:}
		& \begin{tabular}[t]{l}
			Reyna Patricia \\ González Rodríguez
		\end{tabular} \\
		& \\
		\multicolumn{2}{l}{Fecha de realización: \today} \\
		\multicolumn{2}{l}{\universitylocation}
	\end{tabular}
}

% IMPORTACION DEL TEMPLATE
\input{template}

% FORMATO SOLICITADO / USADO EN ACTIVIDADES CORREGIDAS
% La planeación solicita fuente Arial 12. En LaTeX clásico (pdfLaTeX) se usa Helvetica como sustituto:
\usepackage{helvet}
\renewcommand{\familydefault}{\sfdefault}
% Si desea usar Arial exactamente, compile con XeLaTeX o LuaLaTeX y descomente:
% \usepackage{fontspec}
% \setmainfont{Arial}
% Nota: verificar que el tamaño de letra sea 12pt en el documento final.

% PRODUCTOS VISUALES EN TIKZ
% Usar TikZ para construir productos solicitados cuando sea posible:
% cuadro comparativo, tabla, mapa conceptual, mapa mental, linea de tiempo,
% diagrama de flujo, mapa de cajas, redes conceptuales, esquemas, matrices,
% SQA, Venn, procesos, cronologias, arboles de decision y organizadores.
\usepackage{tikz}
\usetikzlibrary{
	arrows.meta,
	positioning,
	calc,
	matrix,
	fit,
	shapes.geometric,
	shadows.blur
}

% CITAS EN FORMATO AUTOR-ANIO, COMPATIBLE CON APA 7 EN EL CUERPO DEL TEXTO
% Usar:
% - Cita parentetica: \citep{clave}
% - Cita narrativa: \citet{clave}
% - Varias fuentes: \citep{clave1,clave2}
% - Parafrasis con pagina si conviene: \citep[p.~00]{clave}
% - Cita textual breve: "texto exacto" \citep[p.~00]{clave}
%
% Regla practica:
% - Si la consigna pide "con tus propias palabras", usar parafrasis con cita.
% - Si se toma una frase exacta, usar comillas y pagina.
% - En definiciones de glosario, citar inmediatamente despues de la definicion
%   base y desarrollar el analisis propio en las oraciones siguientes.
\setcitestyle{authoryear,open={(},close={)}}

% ==============================================================================
% AYUDAS DE REDACCION
% Quitar o comentar \pendiente cuando el documento final este terminado.
% ==============================================================================
\newcommand{\pendiente}[1]{\textcolor{red}{[PENDIENTE: #1]}}

% ==============================================================================
% MARCA DE AGUA EN PORTADA
% - Se aplica solo a la portada porque usa \AddToShipoutPictureBG*.
% - Para apagarla en alguna actividad: \def\coverwatermarkenabled {false}
% - Usar opacidad baja para que no compita con titulo, datos ni logotipo.
% ==============================================================================
\def\coverwatermarkenabled {true}
\def\coverwatermarkimage {img/departamentos/UnADM.pdf}
\def\coverwatermarkopacity {0.16}
\def\coverwatermarkwidth {0.72\paperwidth}
\def\coverwatermarkxshift {0cm}
\def\coverwatermarkyshift {-0.35cm}

\newcommand{\insertcoverwatermark}{%
	\ifthenelse{\equal{\coverwatermarkenabled}{true}}{%
		\AddToShipoutPictureBG*{%
			\begin{tikzpicture}[remember picture,overlay]
				\node[
					opacity=\coverwatermarkopacity,
					inner sep=0pt,
					xshift=\coverwatermarkxshift,
					yshift=\coverwatermarkyshift
				] at (current page.center) {%
					\includegraphics[width=\coverwatermarkwidth]{\coverwatermarkimage}%
				};
			\end{tikzpicture}%
		}%
	}{}%
}

\begin{document}

% PORTADA
\insertcoverwatermark
\templatePortrait

% CONFIGURACION DE PAGINA Y ENCABEZADOS
\templatePagecfg
\onehalfspacing

% ==============================================================================
% RESUMEN / ABSTRACT
% Debe ser breve, no copiar instrucciones.
% Formula recomendada:
% Tema + objetivo + enfoque/metodo + postura/resultado.
% ==============================================================================
\begin{abstractd}
	\textbf{Objetivo:} Analizar la interpretación jurídica y sus métodos (literal, histórica, teleológica y sistemática), incorporar hermenéutica y argumentación, y aplicar estos criterios al estudio de dos tesis de la Suprema Corte de Justicia de la Nación sobre el delito de violación.

	\textbf{Postura de análisis:} La interpretación jurídica no es un trámite técnico; es el punto donde el sistema decide qué derechos protege y qué límites impone al poder punitivo. Si la interpretación se ancla en estereotipos o formalismos, la norma pierde su función de garantía; si se sostiene con hermenéutica y argumentación, puede corregir sesgos y producir decisiones más justas.
\end{abstractd}

% TABLA DE CONTENIDOS - INDICE
\templateIndex

% CONFIGURACIONES FINALES
\templateFinalcfg

% ======================= INICIO DEL DOCUMENTO =======================

% ==============================================================================
% SECCION 1: INTRODUCCION
% Apertura recomendada:
% - Iniciar con una afirmacion fuerte.
% - Traducir la consigna a problema juridico.
% - Mostrar por que el tema importa.
% - Cerrar anticipando el enfoque personal.
%
% Evitar:
% - "En este trabajo hablare de..."
% - "Es muy importante..."
% - Definiciones de diccionario sin interpretacion.
% ==============================================================================
\section{Introducción}

Interpretar una norma no es solo leerla; es definir su sentido operativo dentro del sistema jurídico. En ese acto se decide qué derechos se protegen, qué conductas se consideran legítimas y qué límites se imponen al poder punitivo. Por eso la interpretación funciona como un criterio de control, no como un trámite mecánico.

En materia penal, la forma de interpretar conceptos como violencia o consentimiento determina si el sistema reconoce la vulnerabilidad de la víctima o reproduce estereotipos. Las dos tesis de la SCJN sobre el artículo 175 de los códigos penales de Jalisco y de la Ciudad de México muestran este giro interpretativo al ampliar el entendimiento de la violencia física y al descartar la incapacidad de resistencia como consentimiento \citep{scjnViolenciaFisica2022,scjnIncapacidadResistencia2019}.

La reflexión que sigue parte de los métodos de interpretación, la hermenéutica y la argumentación jurídica, y después aplica esos criterios al análisis de ambos precedentes. La posición de fondo es que interpretar exige justificar con razones y con efectos reales: lo que se decide en una sentencia termina definiendo la protección efectiva de los derechos.

\subsection*{Criterio de análisis}

Para evitar que el trabajo se convierta en un resumen de fuentes, el análisis se organiza en dos movimientos. Primero se fija una postura conceptual sobre la interpretación jurídica; después se revisa cada precedente a partir de hechos, problema jurídico, norma, criterios interpretativos, argumentos y consecuencia. Esa estructura ayuda a mantener una redacción personal: no se repite la tesis judicial, sino que se explica qué controla, qué corrige y qué riesgos interpretativos deja abiertos.

% ==============================================================================
% SECCION 2: DESARROLLO / ANALISIS
% Organizar segun el producto solicitado:
% - Mapa conceptual: concepto central, conceptos principales, conceptos secundarios,
%   relaciones y conclusion.
% - Cuadro comparativo: categorias, criterios, autores, fundamentos, principios,
%   caracteristicas, ejemplos, limites y valoracion critica.
% - Glosario: termino, definicion operativa, fuente, ejemplo juridico y comentario.
% - Estudio de caso: hechos, problema juridico, normas/principios, analisis,
%   postura y solucion.
%
% Regla derivada de Actividad 3:
% - Cuando la actividad pida glosario y analisis de caso, no basta con enumerar
%   conceptos. Ubicar primero la funcion de los conceptos y despues explicar cada
%   uno con definicion base + cita + comentario propio.
% - En la tabla del caso, no repetir la definicion del glosario: aplicar el
%   concepto a hechos, norma, criterio judicial, consecuencia o tension moral.
%
% Regla de parrafo:
% Entrada: idea/concepto.
% Proceso: como funciona.
% Salida: por que importa o que consecuencia produce.
% ==============================================================================
\section{Reflexión}

La interpretación jurídica es el proceso mediante el cual se atribuye sentido a las normas para resolver problemas concretos; no es un acto neutral, sino una decisión que conecta texto, contexto y consecuencias. Por eso la interpretación opera como un filtro de validez práctica: define cómo se protege un derecho y qué conductas se sancionan \citep{hernandezManriquezHermeneutica2019}.

\subsection{Métodos de interpretación}

\begin{itemize}
	\item \textbf{Interpretación literal.} Parte del significado gramatical de la norma. Su ventaja es la seguridad jurídica, pero su límite es el formalismo cuando el texto se vuelve insuficiente para captar el daño real \citep{scjnMemoriaArgumentacion2008}.
	\item \textbf{Interpretación histórica.} Busca la intención del legislador y el contexto en que nació la norma. Ayuda a fijar alcance, pero puede congelar el sentido y dejar fuera transformaciones sociales \citep{hernandezManriquezHermeneutica2019}.
	\item \textbf{Interpretación teleológica.} Atiende al fin de la norma. Permite recuperar su función protectora, aunque requiere argumentar con cuidado para evitar lecturas arbitrarias \citep{scjnMemoriaArgumentacion2008}.
	\item \textbf{Interpretación sistemática.} Integra la norma dentro del sistema jurídico y sus principios. Da coherencia y evita contradicciones, especialmente cuando se conecta con derechos humanos \citep{scjnIncapacidadResistencia2019,scjnViolenciaFisica2022}. En particular, la interpretación sistemática debe articular mandatos constitucionales y normas de protección de víctimas (p. ej. \citep{cpeum2026,LeyGeneralVictimas}), garantizando que la lectura de la norma responda a obligaciones constitucionales y convencionales.
\end{itemize}

La hermenéutica jurídica amplía la mirada al reconocer que toda interpretación tiene un horizonte de sentido: no basta traducir palabras, hay que comprender prácticas, estructuras de poder y efectos sobre las personas \citep{hernandezManriquezHermeneutica2019}. La argumentación jurídica, por su parte, exige justificar por qué un método es preferible a otro en un caso concreto; sin esa justificación, la interpretación pierde legitimidad y se vuelve discrecional \citep{scjnMemoriaArgumentacion2008}. Desde esta perspectiva, interpretar bien no es elegir el método más cómodo, sino el que resuelve el problema con mayor coherencia, control y protección de derechos.

% ==============================================================================
% BLOQUE OPCIONAL: GLOSARIO Y ANALISIS DE CASO
% Usar cuando la planeacion pida "Glosario", "conceptos juridicos fundamentales",
% "analisis de caso" o una tabla concepto/aplicacion practica.
%
% Reglas derivadas de la retroalimentacion:
% - Minimo de conceptos: respetar el numero de la planeacion.
% - Orden alfabetico si se solicita.
% - Cada concepto: definicion propia + cita APA + funcion juridica.
% - No usar citas textuales en todos los conceptos si se piden palabras propias.
% - Cita textual solo si es indispensable; incluir pagina.
% - Antes del glosario, se puede incluir una ubicacion conceptual breve, no una
%   explicacion larga de la consigna.
% - En el caso: identificar al menos los conceptos pedidos y explicar como operan.
% - Conclusion: actos juridicos identificados + relacion con moral/derecho +
%   postura personal.
% ==============================================================================
%
% \section{Glosario de conceptos juridicos fundamentales}
%
% \subsection*{Ubicacion conceptual previa}
% \pendiente{Clasificar brevemente los conceptos por funcion: sistema juridico,
% norma, sujetos, relacion juridica, hechos/actos, moral/derecho. No copiar la
% instruccion de la planeacion.}
%
% \begin{description}
% \item[\textbf{Concepto.}]
% Definicion base en palabras propias \citep{claveFuente}. Funcion del concepto
% dentro del sistema juridico. Ejemplo breve o consecuencia si se interpreta mal.
%
% \item[\textbf{Otro concepto.}]
% Definicion base en palabras propias \citep[p.~00]{claveFuente}. Comentario
% propio: que problema resuelve, que limite tiene y como se aplica.
% \end{description}
%
% \section{Analisis de caso}
%
% \begingroup
% \small
% \setlength{\tabcolsep}{5pt}
% \renewcommand{\arraystretch}{1.35}
% \begin{longtable}{@{}p{0.24\linewidth}p{0.70\linewidth}@{}}
% \caption{Conceptos juridicos aplicados al caso}\label{tab:analisis-caso}\\
% \toprule
% \textbf{Concepto} & \textbf{Aplicacion practica en el caso} \\
% \midrule
% \endfirsthead
% \toprule
% \textbf{Concepto} & \textbf{Aplicacion practica en el caso} \\
% \midrule
% \endhead
% \bottomrule
% \endlastfoot
%
% \textbf{Concepto juridico} &
% Aplicar el concepto al hecho, norma, autoridad, derecho, deber, sancion,
% validez, eficacia o tension moral del caso. No repetir la definicion. \\
% \midrule
%
% \end{longtable}
% \endgroup

% ==============================================================================
% BLOQUE OPCIONAL: CUADRO COMPARATIVO
% Usar cuando la actividad pida cuadro comparativo.
% Comentarios de rubrica:
% - Categorias claras y jerarquizadas.
% - No solo datos: incluir funcion, limites, aplicacion y valoracion critica.
% - Al final del cuadro debe venir conclusion/postura personal si la planeacion lo pide.
% ==============================================================================
%
% \clearpage
% \pagestyle{empty}
% \begin{landscape}
% \begingroup
% \scriptsize
% \setlength{\tabcolsep}{4pt}
% \renewcommand{\arraystretch}{1.35}
% \begin{longtable}{@{}p{0.18\linewidth}p{0.39\linewidth}p{0.39\linewidth}@{}}
% \caption{Cuadro comparativo entre [tema A] y [tema B]}\label{tab:cuadro-comparativo}\\
% \toprule
% \textbf{Criterio} & \textbf{Tema A} & \textbf{Tema B} \\
% \midrule
% \endfirsthead
% \toprule
% \textbf{Criterio} & \textbf{Tema A} & \textbf{Tema B} \\
% \midrule
% \endhead
% \midrule
% \multicolumn{3}{r}{Continua en la siguiente pagina} \\
% \endfoot
% \bottomrule
% \endlastfoot
%
% Concepto &
% [Definicion con funcion, no solo descripcion. Cita APA.] &
% [Definicion con funcion, no solo descripcion. Cita APA.] \\
% \midrule
%
% Fundamento &
% [Base teorica, juridica o filosofica.] &
% [Base teorica, juridica o filosofica.] \\
% \midrule
%
% Principios &
% [Principios ordenados por jerarquia.] &
% [Principios ordenados por jerarquia.] \\
% \midrule
%
% Aplicacion &
% [Ejemplo real o normativo.] &
% [Ejemplo real o normativo.] \\
% \midrule
%
% Valoracion critica &
% [Fortaleza, limite y costo si se aplica mal.] &
% [Fortaleza, limite y costo si se aplica mal.] \\
%
% \end{longtable}
% \endgroup
% \end{landscape}
% \pagestyle{fancy}

% ==============================================================================
% BLOQUE OPCIONAL: MAPA CONCEPTUAL O FIGURA
% Usar si la actividad pide mapa conceptual, esquema o evidencia visual.
% Si se usa herramienta externa, insertar captura y, si procede, enlace.
%
% Reglas por retroalimentacion:
% - Si la planeacion pide mapa conceptual, el producto principal debe ser mapa,
%   no un texto expositivo disfrazado.
% - Incluir concepto central claramente identificado.
% - Incluir conceptos principales y secundarios.
% - Usar palabras de enlace en las flechas o conexiones.
% - Mostrar jerarquia logica: no repetir conceptos en el mismo nivel sin relacion.
% - Agregar conclusion breve solo si la planeacion lo pide o ayuda a cerrar.
% ==============================================================================
%
% \clearpage
% \begin{figure}[H]
% 	\centering
% 	\includegraphics[width=0.95\linewidth]{img/nombre-del-mapa.png}
% 	\caption{Mapa conceptual de [tema]}
% 	\label{fig:mapa-conceptual}
% \end{figure}

% ==============================================================================
% SECCION 3: POSTURA / ANALISIS PERSONAL
% Esta seccion evita la observacion de "informacion general" o "no se detecta
% autenticidad". Debe contestar:
% - Que posicion se sostiene.
% - Por que esa posicion es razonable.
% - Que evidencia/fuente la respalda.
% - Como se aplica al sistema juridico mexicano.
% ==============================================================================
\section{Análisis de caso}

\subsection{Criterio 1: violencia física en el delito de violación (Jalisco)}

\textbf{Hechos.} En el caso analizado, la imputación por violación se sustentó en que la víctima perdió la consciencia tras ingerir alcohol y el imputado aprovechó esa situación para imponer la cópula. El juez de control dictó un auto de no vinculación a proceso al considerar que no se acreditaba violencia física, pero la instancia revisora revocó esa decisión. El criterio sostiene que la violencia física no exige necesariamente fuerza ni resistencia explícita \citep{scjnViolenciaFisica2022}.

\textbf{Norma aplicable.} Artículo 175 del Código Penal para el Estado de Jalisco, que tipifica el delito de violación y exige violencia física como medio de comisión \citep{scjnViolenciaFisica2022}.

\textbf{Problema jurídico.} Determinar si la violencia física debe limitarse a fuerza corporal directa y resistencia visible, o si también comprende el aprovechamiento de una condición de vulnerabilidad que impide a la víctima reaccionar. El punto crítico es que una lectura estrecha puede dejar fuera conductas que afectan la libertad sexual aunque no presenten la imagen tradicional de agresión.

\textbf{Interpretación literal.} La lectura gramatical podría entender la violencia física como fuerza corporal directa y resistencia de la víctima; bajo esta lectura, el caso quedaría fuera si no se acredita oposición física.

\textbf{Interpretación histórica.} Las concepciones tradicionales del delito exigían fuerza y resistencia como signos visibles de violencia; este antecedente explica por qué el juez de control adoptó un criterio restrictivo.

\textbf{Interpretación teleológica.} El fin de la norma es proteger la libertad sexual y la integridad personal; si el agresor se aprovecha de un estado de vulnerabilidad, se actualiza la violencia física como medio de imposición, aun sin golpes o resistencia.

\textbf{Interpretación sistemática.} La norma debe leerse de forma coherente con el sistema de derechos humanos y con la protección de víctimas; por ello, la violencia física se entiende de manera amplia cuando existe control y aprovechamiento de la vulnerabilidad.

\textbf{Postura y consecuencia.} Desde mi lectura, el criterio es correcto porque desplaza el análisis desde la reacción de la víctima hacia la conducta del agresor. Esa diferencia no es menor: si el sistema exige resistencia visible, termina castigando a la víctima por no haber reaccionado; si observa el control ejercido sobre ella, la interpretación recupera la finalidad protectora de la norma. El límite está en exigir prueba suficiente sobre la condición de vulnerabilidad y sobre el aprovechamiento de esa condición.

\textbf{Argumentos a favor.}
\begin{enumerate}
	\item Actualiza el concepto de violencia física y evita estereotipos que exigen resistencia para reconocer la agresión.
	\item Fortalece la tutela de la libertad sexual al reconocer que la vulnerabilidad también puede constituir un medio de imposición.
	\item Mejora el acceso a la justicia al eliminar criterios probatorios que suelen invisibilizar la violencia en contextos de incapacidad momentánea.
\end{enumerate}

\textbf{Argumentos en contra.}
\begin{enumerate}
	\item Puede ampliar en exceso el concepto de violencia física si no se delimitan con claridad las condiciones de vulnerabilidad.
	\item Existe el riesgo de que la interpretación judicial invada el ámbito legislativo al reconfigurar un elemento típico sin reforma legal expresa.
	\item La indeterminación del concepto puede afectar la seguridad jurídica si no se establecen parámetros probatorios consistentes.
\end{enumerate}

\subsection{Criterio 2: incapacidad de resistencia y consentimiento en la violación equiparada (CDMX)}

\textbf{Hechos.} El criterio aborda la violación equiparada cuando la víctima no puede resistir o no comprende lo que ocurre. La tesis sostiene que esas condiciones no pueden considerarse consentimiento, porque el consentimiento exige capacidad real de decisión \citep{scjnIncapacidadResistencia2019}.

\textbf{Norma aplicable.} Artículo 175 del Código Penal del Distrito Federal, aplicable para la Ciudad de México, que tipifica la violación equiparada en casos de incapacidad de resistencia o ausencia de comprensión \citep{scjnIncapacidadResistencia2019}.

\textbf{Problema jurídico.} Definir si la ausencia de resistencia puede interpretarse como consentimiento, o si el consentimiento requiere capacidad real de comprender, decidir y expresar una voluntad libre. La consecuencia jurídica es directa: confundir pasividad con consentimiento reduce la protección penal y desplaza indebidamente la carga hacia la víctima.

\textbf{Interpretación literal.} El texto legal distingue entre violación básica y equiparada; en esta última no se exige violencia física o moral, sino incapacidad de resistencia o comprensión.

\textbf{Interpretación histórica.} El legislador incorporó la figura para proteger a personas en condiciones de vulnerabilidad que no pueden oponerse o comprender, ampliando la tutela penal frente a agresiones sexuales.

\textbf{Interpretación teleológica.} El propósito es proteger la libertad y autonomía sexual, evitando que la pasividad o la incapacidad se interpreten como aceptación válida.

\textbf{Interpretación sistemática.} La norma se integra con principios de dignidad, igualdad y protección reforzada de personas vulnerables; por ello, el consentimiento solo es válido cuando existe capacidad real de decisión.

\textbf{Postura y consecuencia.} Considero que esta tesis fortalece la interpretación constitucional del consentimiento porque no lo trata como una simple ausencia de oposición, sino como una decisión válida. Sin embargo, el análisis debe cuidarse para no convertir la protección en paternalismo: la clave no es negar autonomía, sino verificar si en el caso concreto existía capacidad efectiva para decidir. Ese control evita tanto la impunidad como una aplicación automática de la norma.

\textbf{Argumentos a favor.}
\begin{enumerate}
	\item Distingue entre consentimiento y mera tolerancia, evitando que la incapacidad se convierta en una excusa para la impunidad.
	\item Protege a personas en situación de vulnerabilidad sin exigir resistencia física como criterio probatorio.
	\item Alinea la interpretación con la tutela reforzada de derechos humanos y con la exigencia de consentimiento libre e informado.
\end{enumerate}

\textbf{Argumentos en contra.}
\begin{enumerate}
	\item Si se aplica de forma automática, puede derivar en lecturas paternalistas que desconozcan la autonomía sexual de personas con discapacidad.
	\item La frontera entre incapacidad y consentimiento puede volverse ambigua si no se analizan con rigor las circunstancias concretas del caso.
	\item Un estándar demasiado amplio puede generar incertidumbre jurídica si no se acompaña de criterios probatorios claros.
\end{enumerate}

% ==============================================================================
% SECCION 4: CONCLUSION
% Debe cerrar con fuerza:
% - Sintetizar el hallazgo.
% - Reafirmar postura.
% - Indicar aplicacion real.
% - Evitar repetir la introduccion.
%
% Si la actividad pregunta "cual corriente tiene mayor relevancia y por que",
% responder de forma directa, no dejarlo implicito.
% ==============================================================================
\section{Conclusiones}

La interpretación jurídica no opera en el vacío: define el alcance real de los derechos y la capacidad del sistema para responder ante daños concretos. Los métodos interpretativos, la hermenéutica y la argumentación muestran que el sentido de una norma no se agota en el texto, sino en su función protectora. En los dos criterios analizados, la SCJN mueve el foco desde la resistencia física hacia la protección de la autonomía sexual, corrigiendo estereotipos que históricamente han limitado el reconocimiento de la violencia \citep{scjnViolenciaFisica2022,scjnIncapacidadResistencia2019}.

Mi postura es que estos criterios son compatibles con un enfoque de derechos, pero solo serán legítimos si se aplican con rigor probatorio y argumentos consistentes. Interpretar no es ampliar por ampliar; es justificar por qué una lectura es más coherente con la finalidad de la norma y con el sistema constitucional. Para la práctica jurídica, esto implica asumir que cada decisión interpretativa tiene un costo y una consecuencia: o fortalece la protección de las personas, o reproduce los límites de un formalismo que ya no responde a la realidad.

% ==============================================================================
% REFERENCIAS
% Reglas:
% - Toda fuente citada en el texto debe estar en el .bib.
% - Toda referencia en el .bib debe estar citada en el texto final; quitar
%   fuentes que no se usaron.
% - Toda cita debe tener clave correcta y coincidir con autor/anio visible.
% - Usar fuentes de la planeacion y, si se complementa, que sean confiables.
% - En el cuerpo: \citep{clave}; no usar solo URLs sueltas.
% - Si se trabaja con PDF local, registrar datos bibliograficos completos:
%   autor, anio, titulo, editorial/institucion y URL/DOI si existe.
% - Si la fuente no tiene fecha, usar s. f. con cuidado y preferir confirmar los
%   datos en la portada o ficha bibliografica del documento.
% ==============================================================================
\bibliography{filosofia-del-derecho}

% FIN DEL DOCUMENTO
\end{document}
